\setcounter{section}{1}
\setlength{\parindent}{0pt}  % niente rientro
\setlength{\parskip}{6pt}  
\section{Analisi dei componenti}
Questa sezione è dedicata alla presentazione dell'architettura
adottata in termini di componenti(CMP), definiti
sulla base dei requisiti funzionali descritti precedentemente.
Per rappresentare la loro interconnessione è utilizzato un
diagramma dei componenti, identificando le interfacce fornite
tra i componenti e verso sistemi esterni


\subsection{Definizione dei componenti}
%DESCRIZIONE COMPONENTI PRINCIPALI, SPECIFICANDO GLI INPUT GLI OUTPUT
\subsubsection*{CMP1: Gestione autenticazione}
\textbf{Descrizione:} il componente si occupa della funzionalità
di registrazione di un account e accesso tramite delle credenziali
Comprende la pagina di registrazione, login e la funzionalità di
logout. Permette di identificare coloro che effettuano l'accesso
al sistema. Permettendogli di accedere alle funzionalità e delle
risorse a loro riservate in base al loro ruolo nell'applicazione.
L'utente può selezionare una modalità di accesso alternativo(google, apple)
così da velocizzarne il processo.

\textbf{Interfaccia richiesta - credenziali di accesso:}
le credenziali includono e-mail, password, dati anagrafici ed eventualmente
tecnici del proprio impianto di produzione. Questi sono richiesti in fase di
registrazione. Quando si effettua l'accesso sono richieste solo email e
password.

%\textbf{Interfaccia richiesta - credenziali d'autenticazione di ruolo:}

\textbf{Interfaccia richiesta - autorizzazione autenticazione da Google:}
autorizzazione all’accesso al sistema proveniente dall’SSO di Google che, previo controllo
delle credenziali inserite dall’utente nei sistemi Google, conferma che tali credenziali sono
corrette.

\textbf{Interfaccia richiesta - autorizzazione autenticazione da Apple:}
autorizzazione all’accesso al sistema proveniente dall’SSO di Apple che, previo controllo
delle credenziali inserite dall’utente nei sistemi Google, conferma che tali credenziali sono
corrette.

\textbf{Interfaccia fornita - richiesta autenticazione verso Google:}
indirizzamento dell’utente verso i sistemi SSO Google per l’inserimento delle credenziali Google.

\textbf{Interfaccia fornita - richiesta autenticazione verso Apple:}
indirizzamento dell’utente verso i sistemi SSO Apple per l’inserimento delle credenziali Apple.

\textbf{Interfaccia fornita - richiesta registrazione:}
l’utente richiede di effettuare la registrazione al portale fornendo i propri dati
(nome, cognome, indirizzo, telefono, email, password) i quali vengono salvati nel database.

\textbf{Interfaccia fornita - data di registrazione:}
la data della prima registrazione viene salvata all'interno del database così da ottenere una maggior
storicità dell'account.

\textbf{Interfaccia fornita - richiesta di accesso:}
l’utente richiede di effettuare l'accesso all' applicazione fornendo le credenziali d'accesso.

\textbf{Interfaccia fornita - autenticazione:}
permette a colui che effettua l’accesso di accedere all’applicazione in modo tale che possa
spostarsi agevolmente tra le varie sezioni della stessa utilizzandone i servizi e le funzioni,
permettendo l’accesso ai dati a lui riferiti e la visualizzazione di sezioni a lui concesse.

\subsubsection{CMP2: Gestione account}
\textbf{Descrizione} il componente si occupa delle funzionalità di gestione dell'account, visualizzazione
    e modifica delle proprie informazioni. Comprende la pagina dell'area personale e le
    funzionalità di modifica password, email, numero telefonico o informazioni personali, oltre a quella di
    cancellazione del profilo. Donando il pieno controllo dei propri dati e del proprio account
    all'utente autenticato.
\subsection{Diagramma dei componenti}
%RAPPRESENTAZIONE GRAFICA DELLE INTERCONNESSIONE FRA I COMPONENTI